\section{结语}

国产算力推理引擎的发展，正在从早期的单点适配转向面向真实服务负载的系统化竞争。本文的讨论表明，评价一条技术路线是否成立，不能只看单个算子、单个模型或单次 benchmark 上的峰值结果，更要看它能否在统一架构下同时处理后端异构、复杂任务负载与上游快速演进三方面压力，并最终在国产芯片、国产基础软件栈和本土部署环境中形成可持续的工程闭环。这一判断在近两年的模型与系统演进中已经表现得相当清楚。模型侧，DeepSeek-V2/V3 所代表的 MoE 与 MLA 路线、Qwen2-VL 与 InternVL2 所代表的多图文档理解链路，以及 GLM-4-Voice、MiniCPM-o 等语音交互模型，都在持续扩大推理系统需要管理的状态空间和阶段异质性\citep{liu2024deepseekv2,deepseekv3,wang2024qwen2vl,chen2024internvl2,zeng2024glm4voice,cui2026minicpmo45}。系统侧，从 Orca、Sarathi-Serve 到 DistServe、Splitwise，再到 XGrammar、Llumnix、ElasticMM，研究重点也已从单纯提高吞吐转向如何组织复杂请求、隔离阶段干扰、控制服务级目标并吸收多模态运行时复杂度\citep{orca2023,sarathi2023,zhong2024distserve,splitwise2024,dong2024xgrammar,sun2024llumnix,liu2025elasticmm}。国产推理引擎如果不能把这些变化转化为稳定的抽象边界与工程机制，就很容易在每一轮模型升级中回到高成本重适配状态。

若从更贴近当前国内推理引擎建设任务的视角压缩全文，可以把本文最终收敛为“一个底座、三条主路径”的判断框架：统一抽象、统一指标和可观测能力构成底座，执行与通信、状态治理、压缩与成本协同构成三条主路径。前者决定系统能否在多平台、多版本和多阶段联调条件下保持可组合、可比较和可复现，后者则分别对应吞吐与时延主路径、长上下文与高并发主路径，以及单位 token 成本与部署经济性主路径。这个框架的意义，不在于替代全文已有的路线比较，而在于把不同路线、不同优化和不同交付目标重新放回同一套系统分析口径中。

就现阶段而言，开源主干后端适配、平台一体化方案与混合式架构三条路线各有适用边界。其中，混合式路线更有希望在生态兼容性与平台特化能力之间取得现实平衡，因此也最可能成为未来一段时期的主流工程选择。但无论采取何种路线，真正决定长期竞争力的仍是同一组能力：平台抽象是否清晰，执行与通信主路径是否足够稳健，状态治理体系是否能支撑长上下文和复杂负载，压缩与成本协同是否真正进入运行时闭环，以及评测与交付闭环能否真实反映业务代价。谁能把这些能力稳定组织起来，谁才更有可能把国产算力有效转化为面向 AGI4S 工作负载的生产力。如果把全文的讨论再压缩一步，本文实际想强调三点。第一，国产推理引擎首先是平台边界问题，而不只是模型部署问题；真正困难的不是“支持某个平台”，而是如何把平台差异转化为可治理的能力契约。第二，随着长上下文、多模态、语音和 Agent 负载成为常态，推理系统越来越像状态密集型运行时，其核心竞争力也越来越取决于请求状态、缓存状态和异常回退状态能否被统一组织，而压缩配置与成本核算也必须作为运行时内生变量进入同一闭环。第三，工程闭环并非附属议题，而是决定技术路线能否形成长期复利的核心条件；没有 merge-safe 的扩展边界、没有近真实负载的评测框架、没有围绕部署和回归的治理工具链，再快的局部优化也很难沉淀为稳定能力。从这个意义上说，国产推理引擎的下一阶段竞争，未必首先体现为某一个 benchmark 上再提升多少百分比，而更可能体现为谁能更快、更稳、更低成本地吸收模型变化并完成生产交付。对产业侧而言，这意味着推理引擎建设要从“围绕单一模型做专项优化”转向“围绕持续演进负载建设系统能力”；对研究侧而言，则意味着后续工作需要更重视长周期维护成本、平台迁移成本、复杂工作负载可复用性以及真实服务指标，而不是停留在孤立的局部加速叙事上。归结起来，本文想强调的结论并不复杂：国产算力推理引擎的核心竞争力，最终取决于统一底座、三条主路径与工程闭环能否被同时建立，并在持续演进中保持稳定。
